\documentclass[11pt,a4paper]{article}

\usepackage[T1]{fontenc}
\usepackage[utf8]{inputenc}
\usepackage[provide=*,french]{babel}
\usepackage[a4paper,margin=2.3cm]{geometry}
\usepackage{hyperref}
\usepackage{enumitem}
\usepackage{booktabs,tabularx}
\usepackage{xcolor}
\usepackage{fancyvrb}

\hypersetup{
  colorlinks=true,
  linkcolor=blue,
  urlcolor=blue,
  pdftitle={myco apps releases},
  pdfauthor={Eddy Boite}
}

\DefineVerbatimEnvironment{CodeBlock}{Verbatim}{fontsize=\small}
\setlist[itemize]{leftmargin=1.2cm}
\setlist[enumerate]{leftmargin=1.2cm}
\setlength{\emergencystretch}{3em}

\title{myco\_apps\_releases}
\author{Eddy Boite}
\date{15 juin 2026}

\begin{document}
\maketitle
\tableofcontents
\newpage

\section{Présentation}

Ce dépôt publie des scripts R prêts à l'emploi issus de projets de mycologie et
d'écologie. Chaque application est fournie avec ses données d'exemple, son script
principal, une documentation d'utilisation et des répertoires de sortie standardisés.

\section{Applications disponibles}

\begin{description}[style=nextline]
  \item[Inventaires fongiques --- Complétude et Représentativité]
  Script ICR version 1.4, dédié à l'analyse de la complétude d'inventaires et à la
  représentativité des observations par site, date et espèce.

  \item[Évaluation du potentiel fongique, intérêt patrimonial et gradient CHEGD]
  Script EPFIP-CHEGD version 1.0, dédié au calcul du potentiel fongique, de l'indice
  patrimonial, du gradient CHEGD et de l'indice de représentativité.
\end{description}

\section{Structure du dépôt}

\begin{CodeBlock}
myco_apps_releases/
|-- data/                         # Donnees d'exemple
|-- docs/                         # Documentation
|-- logs/                         # Journaux d'execution
|-- results/                      # Resultats generes
|   |-- ICR/
|   `-- EPFIP_CHEGD/
|-- scripts/                      # Scripts R principaux
`-- README.md
\end{CodeBlock}

\section{Interface web locale}

Une interface Shiny locale permet de lancer les scripts sans modifier leur code. Elle
fonctionne depuis le poste utilisateur et s'appuie sur les dossiers \texttt{scripts/},
\texttt{data/}, \texttt{logs/} et \texttt{results/}.

\subsection{Lancement}

\begin{itemize}
  \item macOS : double-cliquer sur \texttt{Lancer\_Mac.command}.
  \item Windows : double-cliquer sur \texttt{Lancer\_Windows.bat}.
  \item Linux : exécuter \texttt{./Lancer\_Linux.sh} dans un terminal.
\end{itemize}

\section{Quick Start CHEGD}

\begin{enumerate}
  \item Vérifier que R est installé sur la machine.
  \item Placer le fichier d'entrée dans \texttt{data/} ou sélectionner un fichier depuis l'interface.
  \item Lancer le script CHEGD depuis l'interface ou avec \texttt{Rscript}.
  \item Consulter les résultats dans \texttt{results/EPFIP\_CHEGD/}.
\end{enumerate}

\subsection{Données CHEGD attendues}

\begin{tabularx}{\textwidth}{lX}
\toprule
\textbf{Colonne} & \textbf{Rôle} \\
\midrule
\texttt{Espèces} & Nom scientifique observé. \\
\texttt{Famille} & Famille taxonomique. \\
\texttt{Date} & Date d'observation. \\
\texttt{Nombre d'espèce} & Abondance ou effectif. \\
\texttt{Site} & Site ou pelouse inventoriée. \\
\texttt{Fiabilité détermination} & Niveau de confiance dans la détermination. \\
\bottomrule
\end{tabularx}

\subsection{Sorties CHEGD}

Les sorties principales sont les fichiers de synthèse par site, les tableaux de potentiel
fongique, les indices patrimoniaux, les gradients CHEGD, les indices de représentativité,
les figures au format PNG/PDF et les journaux d'exécution.

\section{Quick Start ICR}

\begin{enumerate}
  \item Vérifier que R est installé et accessible dans le terminal.
  \item Préparer un fichier d'observations avec les colonnes attendues.
  \item Lancer l'analyse ICR depuis l'interface ou avec \texttt{Rscript}.
  \item Consulter les résultats dans \texttt{results/ICR/}.
\end{enumerate}

\subsection{Données ICR attendues}

\begin{tabularx}{\textwidth}{lX}
\toprule
\textbf{Colonne} & \textbf{Rôle} \\
\midrule
\texttt{site} & Identifiant du site inventorié. \\
\texttt{date} & Date de visite. \\
\texttt{visite\_id} & Identifiant unique de visite. \\
\texttt{espece} & Nom d'espèce observée. \\
\texttt{placette} & Colonne optionnelle pour préciser l'unité spatiale. \\
\bottomrule
\end{tabularx}

\section{Installation des dépendances}

Les scripts vérifient et installent les packages R nécessaires lorsque cela est possible.
Une connexion internet est recommandée au premier lancement. En environnement contrôlé,
il est préférable d'installer les dépendances avant la distribution de l'archive.

\section{Interprétation et bonnes pratiques}

\begin{itemize}
  \item Toujours vérifier la qualité taxonomique et la cohérence des dates avant interprétation.
  \item Interpréter les scores à la lumière de l'effort d'échantillonnage et de la saison.
  \item Conserver les logs avec les résultats afin de documenter chaque exécution.
  \item Ne pas comparer deux campagnes si les méthodes ou les colonnes d'entrée ont changé sans traçabilité.
\end{itemize}

\section{Dépannage}

\begin{description}[style=nextline]
  \item[Le script ne trouve pas le fichier d'entrée]
  Vérifier le chemin, le nom de fichier et la présence du fichier dans \texttt{data/}.

  \item[Les packages R ne s'installent pas]
  Vérifier la connexion internet, les droits d'écriture de la bibliothèque R et les messages du log.

  \item[Les résultats sont absents]
  Consulter le dernier fichier de log dans \texttt{logs/} et vérifier que le fichier d'entrée contient les colonnes attendues.
\end{description}

\section{Références}

\begin{itemize}
  \item Wickham, H. (2016). \textit{ggplot2: Elegant Graphics for Data Analysis}. Springer.
  \item Venables, W.N. \& Ripley, B.D. (2002). \textit{Modern Applied Statistics with S}. Springer.
  \item Documentation R et documentation des packages utilisés par les scripts.
\end{itemize}

\end{document}
